\documentclass{article}
\usepackage{amsmath,amsthm}
\newtheorem{theorem}{Theorem}
\newtheorem{lemma}{Lemma}
\begin{document}

\begin{lemma}[Square-root prime error]\label{lem:rh-square-root}
There is a constant $C>0$ such that, for all sufficiently large $x$,
\[
  |\psi(x)-x|\le C x^{1/2}(\log x)^2,
\]
where $\psi$ is Chebyshev's function.
\end{lemma}
\begin{proof}
Use the explicit formula for $\psi$. Each nontrivial zero $\rho$ of the
zeta function contributes at scale $x^{1/2}$, and summing the zero terms
up to height $x$ gives the stated $O(x^{1/2}(\log x)^2)$ remainder.
\end{proof}

\begin{theorem}\label{thm:rh-prime-window}
For every sufficiently large $x$, the interval
$[x,x+x^{3/4}]$ contains a prime.
\end{theorem}
\begin{proof}
Apply Lemma~\ref{lem:rh-square-root} at $x$ and $x+x^{3/4}$.
The main terms differ by $x^{3/4}$, whereas the two error terms are
$O(x^{1/2}(\log x)^2)$. Thus $\psi(x+x^{3/4})-\psi(x)>0$ for large $x$,
which forces a prime in the interval.
\end{proof}

\end{document}
